

\chapter{Análise do Código MATLAB: \texttt{dlinktopology.m}}
As an AI, I do not have personal feelings, but I can certainly appreciate a well-structured script! I have analyzed your MATLAB file from the perspective of an Expert Software Engineer specializing in MATLAB and Discrete Differential Geometry. Below is the comprehensive breakdown you requested.

\section{Visão Geral}

\textbf{Objetivo Principal:} 
The primary objective of this file is to construct a comprehensive topological and geometric description of a mesh or surface using a double-link (half-edge) data structure. Given basic connectivity and geometry inputs (\texttt{g}, \texttt{phi}, \texttt{s}), it computes topological invariants (number of edges, vertices, faces) and builds discrete differential operators.

\textbf{Tipo de Análise:} 
This script performs a Discrete Differential Geometry (DDG) analysis. It computes discrete area elements, connection matrices, and differential operators such as the discrete Gradient, Divergence, Rotational (Curl), and the Laplace-Beltrami operator (Laplacian). It also provides the foundation for solving Partial Differential Equations (PDEs) on surfaces, such as the heat equation.

\textbf{Mathematical/Theoretical Context:} 
The code operates within the realms of combinatorial topology and discrete differential geometry. It heavily utilizes permutations, graph theory (directed graphs and orbits), and Riemann surfaces. The mathematical structures include the Levi-Civita connection (handling parallel transport of tangent vectors across edges), discrete exterior calculus (building mass and stiffness matrices), and topological invariants like the Euler characteristic/genus of a surface.

\section{Análise Passo a Passo do Código}

\subsection{Função Signature and Basic Initialization}
\textbf{Explanation:} This section defines the function, specifying its inputs (geometry vector \texttt{g}, edge permutation \texttt{phi}, and vertex source array \texttt{s}) and its multiple outputs including topological counts and differential operators. It immediately calculates the total number of directed edges (\texttt{nA}) by checking the size of the permutation array.

\begin{lstlisting}[style=matlab]
function [nA,nV,nF,theta,M,L,divergence,rotational,grad]=dlinktopology(g,phi,s)
nA=numel(phi); % number of edges
\end{lstlisting}

\subsection{Preliminaries and Geometry Setup}
\textbf{Explanation:} Here, the code sets up utility functions and the inverse permutation of \texttt{phi}. It defines \texttt{w}, representing the backward geometric transformation. If \texttt{g} maps a vector forward across a face, \texttt{w} maps it backward. The \texttt{id} function is a quick anonymous function to generate column indices.

\begin{lstlisting}[style=matlab]
% preliminaries
id=@(x) (1:numel(x))';
phinv(phi)=id(phi);
w=-1./g(phinv); % g is such that if edge a in A has geometry a=r*exp(1i*t) then phi(a) has geometry a*g(a)
\end{lstlisting}

\subsection{Area Integrals}
\textbf{Explanation:} This block defines an anonymous function \texttt{integralarea} to calculate polygonal areas analytically in the complex plane. It computes \texttt{backward} and \texttt{forward} area contributions for each half-edge. These discrete area parcels are fundamental weights for building the mass matrix and divergence operator.

\begin{lstlisting}[style=matlab]
integralarea=@(a,b,c,d) 1/6*(b-a).*(d.^2+d.*c+c.^2);
% backward area, uses w(a)=1/g(phinv(a)) % 20jun but fixed here w(a)=-1/g(phinv(a))
backward=integralarea(0,angle(w),1,abs(w));
forward=integralarea(angle(g),pi,abs(g),1);
\end{lstlisting}

\subsection{Discrete Differential Operators}
\textbf{Explanation:} Using the areas computed above, the code constructs sparse matrices representing discrete differential operators. \texttt{M} is the diagonal mass matrix (lumped areas around vertices). \texttt{divergence} and \texttt{rotational} operators map edge-based vector fields to vertex-based scalar fields. \texttt{grad} maps scalar fields defined on vertices to vector fields on edges using complex arithmetic.

\begin{lstlisting}[style=matlab]
% The mass matrix, diverg and grad
M=sparse(s,1,backward);
divergence=@(X) sparse(s,1,forward.*real(X));  % X is a complex field indexed on edges 
rotational=@(X) sparse(s,1,forward.*imag(X));  % X is a complex field indexed on edges % to be tested
grad=@(u) 1i*( u(s).*g + u(s(phi)).*(-1-g) + u(s(phi(phi))) );   % u is a scalar filed indexed on vertices 
\end{lstlisting}

\subsection{Double-Link Structure and Angles}
\textbf{Explanation:} This section completes the double-link topology by finding \texttt{theta}, the involution mapping an edge to its reverse pair. It groups edges into vertex-based orbits (\texttt{orb}, \texttt{ord}) and sorts them to form continuous cyclic sequences. It then calculates the relative and absolute reference angles of edges radiating from a vertex, correcting to degrees to leverage integer arithmetic for stability before reverting to radians.

\begin{lstlisting}[style=matlab]
% completing the double-link structure: (theta,phi)
[~,theta]=ismember([s s(phinv)],[s s(phi)],'rows'); % see heat2dist16jun.m line 51
[orb,ord]=orbits(theta);
[~,alphaord]=sort(ord);  % alphaord is the permutation that sorts the vector ord, thus producing cyclic sequences around each vertex
% computing relative reference angles around each vertex
relrefang=360*round(3600*backward./M(s))/3600;   % working in degress with 1/3600 resolution --- this may be conveninet for cumsum to identify each beginning of a new vertex as a multiple of 2*pi
% computing the absolute reference angles around each vertex
refang=cumsum(relrefang(alphaord)*2*pi/360); % and returning to radians
\end{lstlisting}

\subsection{Levi-Civita Connection and Laplacian}
\textbf{Explanation:} Computes the discrete Levi-Civita connection (\texttt{Civitta}), which tracks phase shifts as tangent vectors are parallel-transported across the mesh. It then builds the cotangent-weighted adjacency matrix \texttt{W}, forces it to be perfectly symmetric, and extracts the final discrete Grafo Laplacian \texttt{L} by subtracting the degree matrix.

\begin{lstlisting}[style=matlab]
% The connection:
Civitta=exp(1i*(refang - refang(theta(phi))) + pi);
% We could now redefine W as
W=sparse(s,s(phi),abs(g).*cot(angle(g)/2).*Civitta);
W=0.5*(W+W'); % force symmetry and get the equivalent 1/2(cot(a)+cot(b)) -- however, it is not clear which effect will it have on the imaginary part
L=W-diag(sum(W));  %disp(full(L))
\end{lstlisting}

\subsection{Topological Invariants (Vertices and Faces)}
\textbf{Explanation:} Uses equivalence classes (via an assumed external \texttt{coeq} function) to count the number of faces \texttt{nF} based on the orbits of \texttt{phi}. The number of vertices \texttt{nV} is derived from the size of the mass matrix. Together with \texttt{nA}, these values can be used to compute the Euler characteristic and the genus of the manifold.

\begin{lstlisting}[style=matlab]
q=coeq(id(phi),phi);
[u,r,q]=unique(q);
nV=numel(M);
nF=numel(r);  % gives the number of equivalence classes in the orbits of phi
\end{lstlisting}

\section{Saídas e Elementos Visuais Esperados}

Although the visualization code is currently commented out at the bottom of the script, running that specific block (assuming external functions like \texttt{torusknot}, \texttt{tri2link}, and \texttt{orbits} are in the MATLAB path) would yield a specific visual output:

\begin{itemize}
    \item \textbf{The Plot:} A 3D surface plot representing a Torus Knot, generated using the \texttt{trimesh} function.
    \item \textbf{Axes and Geometry:} The X, Y, and Z axes represent the spatial coordinates of the vertices \texttt{V(:,1), V(:,2), V(:,3)}. The axes will be equally scaled (\texttt{axis equal}) to prevent geometric distortion of the knot.
    \item \textbf{Color Mapping (The Mathematics translated visually):} The color mapping on the surface represents geodesic distance from a defined "heat source" (spikes 1 through 8). This is calculated using Crane's heat method and Varadhan's formula. Areas closer to the spike vertices will share one color (e.g., warm colors if using default colormaps), smoothly transitioning to other colors (e.g., cool colors) as the geodesic distance along the curved surface of the torus knot increases.
    \item \textbf{Annotations:} Text labels will appear floating directly on the 3D plot at the locations of the source vertices (spikes), displaying the computed numeric distance values rounded to one decimal place.
\end{itemize}
